\section{Related Work}
\label{sec:related}

Three converging research threads form the prior art for \sysname.
We summarise each, then state the gap.

\smallskip
\noindent\textbf{Consent management on distributed ledgers.}
MedRec~\cite{azaria2016medrec} pioneered the use of a blockchain
substrate to record patient consent for medical-data access and
permissioning. The system established the durability and auditability
properties we inherit: every consent state transition is an immutable
log entry, and the consent owner retains exclusive control over
modifications. MedRec, however, treats the consent record as the
\emph{endpoint} of the protocol: there is no downstream ML training
whose execution is gated by the consent state, and no mechanism for
demonstrating that withdrawn consent translated into a model update.
Subsequent semantic-web work (GConsent, W3C DPV) formalises consent
vocabularies but operates in an off-chain registry without
enforcement.
\emph{Gap}: consent is recorded but does not influence training.

\smallskip
\noindent\textbf{Machine unlearning.}
The SISA training scheme~\cite{bourtoule2021sisa} pre-shards training
data and caches per-shard intermediate models so that data removal
requires retraining only the affected shard. FedEraser
\cite{liu2021federaser} extends SISA's logic to the federated setting
by replaying aggregation history with a calibration step,
significantly cheaper than full retraining. Gradient-ascent unlearning methods take a single trained model and
perform a small number of negative-learning-rate steps on the data
to be removed. All three families have been validated on standalone
benchmarks but \emph{none specifies how unlearning is triggered, by
whom, or how a deployer demonstrates it actually happened}. The
unlearning literature treats the trigger and verification steps as
out-of-scope.
\emph{Gap}: unlearning is implementable but not verifiable
end-to-end.

\smallskip
\noindent\textbf{Blockchain-based federated learning.}
BlockFL~\cite{kim2020blockchained} records every device-side model
update on-chain to defend against malicious aggregators. BFLC
\cite{li2020bflc} replaces the FL aggregator with a committee-consensus
mechanism. These systems address aggregator-side trust but do not
encode consent at all: a participant who joins the federation
cannot subsequently withdraw without trusting the orchestrator to
respect the request. Membership inference attacks
\cite{shokri2017mia,carlini2022lira,choquettechoo2021labelonly} provide
the canonical operational test for whether a model still memorises a
specific data subject's training contribution, but no published
blockchain-FL system uses MIA scores as the on-chain acceptance
criterion for a state transition.
\emph{Gap}: model updates are auditable but consent semantics and
forgetting verification are absent.

\begin{table*}[t]
\centering
\caption{Feature comparison with the three closest prior systems.
Each cell describes the corresponding system's mechanism for that
dimension; ``Not provided'' marks dimensions the prior system was
not designed to address, not failures of implementation.
Per-operation cost is omitted here; see
Table~\ref{tab:t2}.}
\label{tab:t3}
\footnotesize
\setlength{\tabcolsep}{4pt}
\renewcommand{\arraystretch}{1.05}
\begin{tabular}{p{2.2cm}p{4.3cm}p{3.7cm}p{3.9cm}p{3.0cm}}
\toprule
Dimension & ConsentChain (ours) & MedRec~\cite{azaria2016medrec} &
FedEraser~\cite{liu2021federaser} & BlockFL~\cite{kim2020blockchained} \\
\midrule
Consent representation &
6-state \clsm on a soulbound ERC-721 token; every transition
signed by the data subject's wallet &
Per-record access policy on a permission contract, owned by the
patient &
Not provided; consent is managed off-chain by the deployer &
Not provided; participation is implicit on join \\

Unlearning trigger \& method &
Chain-gated trigger ($\dmax$-bounded liveness,
Property~\ref{prop:liveness}); method-agnostic, accepting any
(pre, post) model-hash pair from SISA-Fed, gradient ascent,
FedEraser, or future methods &
Not applicable: no downstream ML training &
Trigger left to the deployer; method fixed to FedEraser's
history-replay-with-calibration scheme &
Not applicable: no unlearning step \\

Verification of unlearning &
Three independent MIAs (shadow, LiRA, label-only) all below
$\tauMIA$ \emph{and} $t$-of-$n$ ECDSA signatures from a
registered verifier committee (Property~\ref{prop:verifiability}) &
Not applicable &
Self-reported by the data controller; no external check &
Not applicable \\

On-chain footprint \& audit trail &
3 Solidity contracts (registry, oracle, gatekeeper) on either
gas-free PureChain or gas-based Polygon Amoy; every \clsm
transition, UA acceptance, and \texttt{ComplianceViolation}
emitted as an indexed event &
Permission contracts on Ethereum mainnet; public ledger of
permission grants/revocations &
None; entirely off-chain, no immutable record &
Per-round device-side model-update log on Ethereum mainnet;
public ledger of update commitments \\

\bottomrule
\end{tabular}
\end{table*}

\smallskip
\noindent\textbf{Positioning.}
Table~\ref{tab:t3} contrasts \sysname with the three closest prior
systems across the four dimensions implicated by the gaps above:
consent representation, unlearning trigger and method, verification
of unlearning, and on-chain footprint. Each prior system covers at
most one dimension natively; \sysname is the first design we are
aware of to address all four in a single chain-enforced protocol,
and the first to do so at gas-free per-operation cost on a public
EVM-compatible substrate. We adopt MedRec's on-chain consent
registry pattern (vertex A), extend it with the \clsm safety and
liveness properties; we accept any of the three established
unlearning paradigms as a black-box (vertex B); and we make the
verification step (vertex C) the chain's gating condition for a
consent state transition. The resulting pattern (consent change
$\rightarrow$ unlearning execution $\rightarrow$ three-MIA
verification $\rightarrow$ on-chain attestation $\rightarrow$
consent state finalisation) is the missing edge.

The choice of \emph{which} blockchain matters at IoT scale.
Permissionless gas-based chains (Ethereum, Polygon) impose recurring
per-operation costs that scale linearly with device count and consent
operation rate. Section~\ref{sec:evaluation} quantifies this against
PureChain~\cite{kim2025purechain}, a public, gas-free,
EVM-compatible chain with PoA$^2$ consensus and
Merkle-proof-backed state. PureChain is not the
contribution of this paper; rather, our contribution is showing that
the protocol is operationally identical on both chain types, while
the economic model is fundamentally different.
